% !TeX root = ../papers/nc-p3.tex
\section{二元障害の二方向が重なる十四次因子}
\label{sec:binary-reverse-orientation}
\begin{theorem}\label{thm:binary-reverse-orientation}
階数2の二元障害の行列の核を $k\ell\subset P_2$ とし、
$Y\mathbin{\lrcorner}\ell=T\mathbin{\lrcorner}c$ で四次式 $c$ を定める。
$Y$ の代表による $c$ の変更は $\ell P_2$ の加算である。
\begin{enumerate}
\item $\ell\mid c$ なら点グラフは正次元であり、逆も成立する。
\item $\gcd(\ell,c)=1$ なら有限窓の二元商は指定の向きにしか存在しない。
\item $\gcd(\ell,c)=d$ が一次なら有限窓の二元商は両方向に存在する。
固定した $T$ の $\mathbb P^5$ 上、この両向きの軌跡の閉包は被約既約な次数14の超曲面である。
\end{enumerate}
両向きのときは以下の線型変換が逆向きの十次・八次パラメータを計算する。
二組が異なる向きで提示された場合の添字保持同型は、一方をこの変換で戻して
定理\ref{thm:binary-apolar-rank-test}の同じ向きの軌道検査を行うことと同値である。
これら非正則候補の窓の最小周期は2で、両方向が許され逆向きのパラメータが同じ軌道に戻る場合に限り1となる。
両向き因子上でも一般最小周期は2である。
\end{theorem}
\begin{proof}[反対四次曲線と逆向きの三平面]
$T\mathbin{\lrcorner}P_4\subset P_6^*$ の零化空間は $I_6$ で、対応する写像は単射である。
従って行列の核条件がちょうど $c$ の存在と一意性を与える。
$\ell\mid c$ なら $Y$ は極性空間による変更後 $\ell P_6$ を消すので第二割線の射影にある。
逆に第二割線にあれば、動かした代表はある二次式の倍数を消し、その二次式は核の一意性から $\ell$ である。

$U=P_2\oplus kz$ 側の反対核は、三つの $P_2$ 座標の二次錐と、
$z\ell-c$ の二つの式を持つ。ここで $c$ は $\operatorname{Sym}^2(P_2)$ への任意の持上げである。
二式は反対積の式から消え、階数8なので核を尽くす。
錐の頂点以外を、母線の二元座標 $t$ から読むと、共通零点は
$[\ell(t)\operatorname{ev}^{(2)}_t:c(t)]$ の閉包と、共通根の上の母線とである。
$\ell,c$ が互いに素なら上記パラメータは基点を持たず、二次曲線への射影から生成的一対一である。
従って整な四次曲線を切り、逆向きの商に必要なねじれ三次曲線を含まない。

共通因子が一次なら $\ell=dl$, $c=dc_3$, $(l,c_3)=1$ とする。
二つの端を逆にして $dP_2\subset Q$ を新しい三平面、
$J:U\simeq P_3$, $J(p,a)=lp+a c_3$ を新しい全四文字とする。
新しい積の値を元の $P_5$ に移す線型写像 $I$ を
\[
         I(lh)=dh\quad(h\in P_4),\qquad I(c_3 r)=L(dr)\quad(r\in P_2)
\]
で定める。二空間の交わりは $lc_3P_1$ であり、
$Y\mathbin{\lrcorner}(dl)=T\mathbin{\lrcorner}(dc_3)$ から $L(dl v)=dc_3v$ ($v\in P_1$)。
従って定義は一致し全空間に及ぶ。像は $dP_4$ を含む。
もしこれだけなら $d$ の根での評価を $B$ に合成すると $dP_2\subset Q$ が核に入るため純関係となる。
有限点Frobenius窓の純関係の不在に反するので $I$ は可逆。
これで $B(u,dr)=I(J(u)r)$ となり逆の商が得られた。
新たなGram行列は $I^tGI$、第四文字 $q_0\notin dP_2$ の積は
$j\mapsto I^{-1} B(J^{-1}j,q_0)$ であり、その積にGram行列を掛ければ逆向きの八次式が戻る。
前定理の一意性により二方向を比較する同型と周期の判定が従う。
\end{proof}
\begin{proof}[次数と既約性]
$\ell$ は行列の二次小行列で与えられ、$c$ の係数は $Y$ と $\ell$ に双線型である。
従って $\operatorname{Res}(\ell,c)$ は $\bar Y$ の斉次次数 $4\cdot2+2\cdot3=14$ の式である。
代表を変えても終結式は不変である。階数1の集合は余次元2なので因子にはならない。
$[\ell]\in\mathbb P(P_2)$ を固定し $Y(\ell I_6)=0$ とすると、
$[\bar Y]$ の選択は射影三空間である。$c\bmod\ell P_2$ の二係数はその射影化前で全射の線型像である。
実際 $Y\mapsto Y\mathbin{\lrcorner}\ell$ は $P_8^*\to P_6^*$ の全射なので、
希望する $T\mathbin{\lrcorner}c$ を指定できる。階数2の選択はこの三空間内で稠密である
（二行が $P_2/k\ell$ の独立な二汎関数となればよい）。
$\ell$ の二つの相異なる根の一方を選びそこで $c$ が零となることは一つの線型条件で、
他方で零とする条件とは独立である。順序付きの根対を動かしたこの超平面束は既約である。
各ファイバーの階数1条件は二行二列の行列式の整な二次曲面なので、この超平面を含まない。
二重根の底でも $c$ の値が零となるのはなお真の超平面で、別成分を増やさない。
一般点ではちょうど一根の単純な消失であり次数に重複はない。

\path{computations/verify_binary_reverse_orientation.py}では $d=x$, $l=x-1$, $c_3=x^3+2$ の例で
両表示の積を直接照合する。新旧十次形式を通常の二項係数付き表示 $F$ とし、
$A=(F,F)_{10}$, $B=(F,F)_8$ に対する絶対不変量 $(B,B)_4/A^2$ を比較すると値が異なる。
従って既約な両向き軌跡の一般点は一周期でない。
\end{proof}
